%=============================================================
% Capitolo 4 — Metodi Iterativi per Sistemi Lineari
%=============================================================
\chapter{Metodi Iterativi per Sistemi Lineari}
\label{ch:iterativi}

\section{Motivazione}

I metodi diretti richiedono $\mathcal{O}(n^3)$ operazioni e $\mathcal{O}(n^2)$ di
memoria: per sistemi di grandi dimensioni (e.g.\ $n > 10^6$, tipici in data science e
simulazione) questo è proibitivo. I \textbf{metodi iterativi} costruiscono una
successione $\{x^{(k)}\}$ che converge alla soluzione $x^*$ di $Ax = b$, e spesso
richiedono solo $\mathcal{O}(n)$ operazioni per iterazione se $A$ è sparsa.

%------------------------------------------------------------
\section{Framework generale: splitting della matrice}
\label{sec:splitting}
%------------------------------------------------------------

L'idea comune è \textbf{splittare} la matrice $A$ come:
\begin{equation}
  A = M - N,
  \label{eq:splitting}
\end{equation}
dove $M$ è una matrice scelta in modo che $Mx = c$ sia facile da risolvere.
Sostituendo in $Ax = b$:
\begin{equation}
  Mx = Nx + b \implies x = M^{-1}Nx + M^{-1}b.
  \label{eq:iter_gen}
\end{equation}

Questo definisce la \textbf{iterazione stazionaria}:
\begin{equation}
  x^{(k+1)} = \underbrace{M^{-1}N}_{B}\, x^{(k)} + \underbrace{M^{-1}b}_{c},
  \label{eq:iter_staz}
\end{equation}
dove $B = M^{-1}N = M^{-1}(M - A) = I - M^{-1}A$ è la \textbf{matrice di iterazione}.

Decomponiamo $A$ nei suoi componenti diagonale, triangolare inferiore e superiore:
\begin{equation}
  A = D + L + U,
  \label{eq:DLU}
\end{equation}
dove:
\begin{itemize}
  \item $D = \operatorname{diag}(a_{11}, \ldots, a_{nn})$ è la parte diagonale;
  \item $L$ è la parte strettamente triangolare inferiore ($l_{ij} = a_{ij}$ per $i > j$);
  \item $U$ è la parte strettamente triangolare superiore ($u_{ij} = a_{ij}$ per $i < j$).
\end{itemize}

%------------------------------------------------------------
\section{Metodo di Jacobi}
\label{sec:jacobi}
%------------------------------------------------------------

\begin{definition}[Metodo di Jacobi]
Lo splitting di Jacobi è $M = D$, $N = -(L + U)$. L'iterazione è:
\begin{equation}
  x^{(k+1)} = D^{-1}\bigl(-(L+U)x^{(k)} + b\bigr),
  \label{eq:jacobi_mat}
\end{equation}
ovvero, per ogni componente $i$:
\begin{equation}
  x_i^{(k+1)} = \frac{1}{a_{ii}}\left(b_i - \sum_{\substack{j=1 \\ j \neq i}}^n
  a_{ij}\, x_j^{(k)}\right), \quad i = 1, \ldots, n.
  \label{eq:jacobi_comp}
\end{equation}
\end{definition}

\begin{remark}
Ogni componente $x_i^{(k+1)}$ si calcola usando \emph{solo} le componenti
dell'iterata precedente $x^{(k)}$: il metodo è \textbf{parallelizzabile}.
\end{remark}

La matrice di iterazione di Jacobi è:
\begin{equation}
  B_J = -D^{-1}(L + U) = I - D^{-1}A.
  \label{eq:B_jacobi}
\end{equation}

%------------------------------------------------------------
\section{Metodo di Gauss-Seidel}
\label{sec:gauss_seidel}
%------------------------------------------------------------

\begin{definition}[Metodo di Gauss-Seidel]
Lo splitting di Gauss-Seidel è $M = D + L$, $N = -U$. L'iterazione è:
\begin{equation}
  (D + L)\,x^{(k+1)} = -U\,x^{(k)} + b,
  \label{eq:gs_mat}
\end{equation}
ovvero, per ogni componente $i$:
\begin{equation}
  x_i^{(k+1)} = \frac{1}{a_{ii}}\left(b_i
  - \sum_{j=1}^{i-1} a_{ij}\, x_j^{(k+1)}
  - \sum_{j=i+1}^{n} a_{ij}\, x_j^{(k)}\right), \quad i = 1, \ldots, n.
  \label{eq:gs_comp}
\end{equation}
\end{definition}

\begin{remark}[Spostamenti successivi]
A differenza di Jacobi, Gauss-Seidel usa immediatamente i valori $x_j^{(k+1)}$ già
calcolati (per $j < i$). Questo accelera tipicamente la convergenza, ma rende il
metodo \textbf{non parallelizzabile} nella forma standard.
\end{remark}

La matrice di iterazione di Gauss-Seidel è:
\begin{equation}
  B_{GS} = -(D + L)^{-1}U.
  \label{eq:B_gs}
\end{equation}

%------------------------------------------------------------
\section{Criteri di convergenza}
\label{sec:convergenza}
%------------------------------------------------------------

\subsection{Condizione necessaria e sufficiente}

\begin{theorem}[Convergenza dei metodi iterativi stazionari]
\label{thm:conv_iter}
La successione $x^{(k+1)} = Bx^{(k)} + c$ converge alla soluzione di $Ax = b$
per ogni $x^{(0)}$ se e solo se:
\begin{equation}
  \rho(B) < 1,
  \label{eq:conv_cond}
\end{equation}
dove $\rho(B)$ è il raggio spettrale della matrice di iterazione $B$.
\end{theorem}

Più piccolo è $\rho(B)$, più rapida è la convergenza.

\subsection{Condizione sufficiente: norma}

\begin{corollary}
Se esiste una norma matriciale $\|\cdot\|$ tale che $\|B\| < 1$, allora il metodo
converge.
\end{corollary}

\subsection{Teorema per matrici a dominanza diagonale}

\begin{definition}[Dominanza diagonale stretta per righe]
$A \in \R^{n \times n}$ è \emph{strettamente a dominanza diagonale per righe} se:
\begin{equation}
  |a_{ii}| > \sum_{\substack{j=1 \\ j \neq i}}^n |a_{ij}|, \quad i = 1, \ldots, n.
  \label{eq:dom_diag}
\end{equation}
\end{definition}

\begin{theorem}[Convergenza con dominanza diagonale]
\label{thm:dom_diag}
Se $A$ è strettamente a dominanza diagonale per righe, allora i metodi di Jacobi e
di Gauss-Seidel convergono entrambi.
\end{theorem}

\begin{proof}[Dimostrazione (Jacobi)]
Per la matrice di iterazione $B_J = -D^{-1}(L+U)$:
\begin{equation}
  \|B_J\|_\infty = \max_i \sum_{j \neq i} \frac{|a_{ij}|}{|a_{ii}|} < 1,
\end{equation}
grazie alla dominanza diagonale stretta. Per il corollario precedente, Jacobi converge.
\end{proof}

\begin{example}[Sistema a dominanza diagonale]
Consideriamo il sistema
\[
  \begin{pmatrix} 10 & -1 & 2 \\ -1 & 11 & -1 \\ 2 & -1 & 10 \end{pmatrix}
  \begin{pmatrix} x_1 \\ x_2 \\ x_3 \end{pmatrix}
  =
  \begin{pmatrix} 6 \\ 25 \\ -11 \end{pmatrix}.
\]
La matrice è strettamente a dominanza diagonale ($10 > 1+2$, $11 > 1+1$, $10 > 2+1$),
quindi Jacobi e Gauss-Seidel convergono entrambi.
\end{example}

%------------------------------------------------------------
\section{Velocità di convergenza e criterio di arresto}
\label{sec:stop}
%------------------------------------------------------------

L'errore al passo $k$ si propaga come:
\begin{equation}
  e^{(k)} = x^{(k)} - x^* = B^k e^{(0)},
\end{equation}
quindi $\|e^{(k)}\| \leq \|B\|^k \|e^{(0)}\|$.

Il criterio di arresto più comune è sul \textbf{residuo relativo}:
\begin{equation}
  \frac{\|b - Ax^{(k)}\|}{\|b\|} < \text{tol},
  \label{eq:stop}
\end{equation}
dove \texttt{tol} è la tolleranza desiderata (e.g.\ $10^{-8}$).

\begin{remark}[Gauss-Seidel vs.\ Jacobi]
In generale $\rho(B_{GS}) \leq \rho(B_J)^2$ (per matrici definite positive
simmetriche): Gauss-Seidel converge circa il doppio più velocemente di Jacobi in
termini di iterazioni. Tuttavia Jacobi è più adatto alla parallelizzazione.
\end{remark}
